\documentclass[10pt,a4paper]{amsart}
\usepackage{amsmath,amssymb,amsthm,graphicx}
% Retained minimal preamble for the Stacks Project chapter more-algebra.tex (More on Algebra).
% The chapter file pinned-original-source.tex carries only \input{preamble} on its first line, so no
% preamble text lives inside the chapter itself.  The Stacks Project uses one single shared preamble for
% every chapter; the snapshot stages it identically as sources/stacks-chow/preamble.tex and
% sources/stacks-etale-cohomology/preamble.tex.  The block below is copied mechanically, line for line,
% from sources/stacks-chow/preamble.tex lines 14-18 (the Xy-pic setup) and 154-217 (the native theorem
% declarations together with the mathematical macros).  Nothing was added, removed or redefined; the
% wrapper class and the amsmath/amssymb/amsthm/graphicx base packages are supplied by the assembler.
% Chapter: more-algebra.tex -- More on Algebra.
% For commutative diagrams we use Xy-pic
\usepackage[all]{xy}

% We use 2cell for 2-commutative diagrams.
\xyoption{2cell}

% Theorem environments.
%
\theoremstyle{plain}
\newtheorem{theorem}[subsection]{Theorem}
\newtheorem{proposition}[subsection]{Proposition}
\newtheorem{lemma}[subsection]{Lemma}

\theoremstyle{definition}
\newtheorem{definition}[subsection]{Definition}
\newtheorem{example}[subsection]{Example}
\newtheorem{exercise}[subsection]{Exercise}
\newtheorem{situation}[subsection]{Situation}

\theoremstyle{remark}
\newtheorem{remark}[subsection]{Remark}
\newtheorem{remarks}[subsection]{Remarks}

\numberwithin{equation}{subsection}

% Macros
%
\def\lim{\mathop{\mathrm{lim}}\nolimits}
\def\colim{\mathop{\mathrm{colim}}\nolimits}
\def\Spec{\mathop{\mathrm{Spec}}}
\def\Hom{\mathop{\mathrm{Hom}}\nolimits}
\def\Ext{\mathop{\mathrm{Ext}}\nolimits}
\def\SheafHom{\mathop{\mathcal{H}\!\mathit{om}}\nolimits}
\def\SheafExt{\mathop{\mathcal{E}\!\mathit{xt}}\nolimits}
\def\Sch{\mathit{Sch}}
\def\Mor{\mathop{\mathrm{Mor}}\nolimits}
\def\Ob{\mathop{\mathrm{Ob}}\nolimits}
\def\Sh{\mathop{\mathit{Sh}}\nolimits}
\def\NL{\mathop{N\!L}\nolimits}
\def\CH{\mathop{\mathrm{CH}}\nolimits}
\def\proetale{{pro\text{-}\acute{e}tale}}
\def\etale{{\acute{e}tale}}
\def\QCoh{\mathit{QCoh}}
\def\Ker{\mathop{\mathrm{Ker}}}
\def\Im{\mathop{\mathrm{Im}}}
\def\Coker{\mathop{\mathrm{Coker}}}
\def\Coim{\mathop{\mathrm{Coim}}}

% Boxtimes
%
\DeclareMathSymbol{\boxtimes}{\mathbin}{AMSa}{"02}

%
% Macros for moduli stacks/spaces
%
\def\QCohstack{\mathcal{QC}\!\mathit{oh}}
\def\Cohstack{\mathcal{C}\!\mathit{oh}}
\def\Spacesstack{\mathcal{S}\!\mathit{paces}}
\def\Quotfunctor{\mathrm{Quot}}
\def\Hilbfunctor{\mathrm{Hilb}}
\def\Curvesstack{\mathcal{C}\!\mathit{urves}}
\def\Polarizedstack{\mathcal{P}\!\mathit{olarized}}
\def\Complexesstack{\mathcal{C}\!\mathit{omplexes}}
% \Pic is the operator that assigns to X its picard group, usage \Pic(X)
% \Picardstack_{X/B} denotes the Picard stack of X over B
% \Picardfunctor_{X/B} denotes the Picard functor of X over B
\def\Pic{\mathop{\mathrm{Pic}}\nolimits}
\def\Picardstack{\mathcal{P}\!\mathit{ic}}
\def\Picardfunctor{\mathrm{Pic}}
\def\Deformationcategory{\mathcal{D}\!\mathit{ef}}

% --- appended 2026-10-12 (Chief): the project own definitions of the reference/slogan/history
% display environments, copied mechanically from the shared Stacks preamble lines 8-12
% (\usepackage{verbatim} plus the three \newenvironment lines, which the chapter preamble keeps
% verbatim from the project preamble). Nothing added, removed or redefined: these environments
% carry no mathematical content and the project itself typesets them as comments. Without these
% three lines a counted statement whose own text uses them fails with Environment slogan undefined.
\usepackage{verbatim}
\newenvironment{reference}{\comment}{\endcomment}
\newenvironment{slogan}{\comment}{\endcomment}
\newenvironment{history}{\comment}{\endcomment}
\usepackage{cleveref}
\crefname{originaldefinition}{original definition}{original definitions}
\Crefname{originaldefinition}{Original definition}{Original definitions}
\makeatletter
\newlabel{definition-types-of-extensions}{{33041}{source}{}{sourcepointer-1}{}}
\newlabel{definition-types-of-extensions@cref}{{[originaldefinition][0][]33041}{[1][0][]source}{}{}{}}
\makeatother
\crefname{originallemma}{original lemma}{original lemmas}
\Crefname{originallemma}{Original lemma}{Original lemmas}
\makeatletter
\newlabel{fields-lemma-subfields-kummer}{{0EXP}{source}{}{sourcepointer-2}{}}
\newlabel{fields-lemma-subfields-kummer@cref}{{[originallemma][0][]0EXP}{[1][0][]source}{}{}{}}
\makeatother
\title{Adjoining an n-th root of a uniformizer gives a totally ramified degree n extension whose integral closure is A[root]}
\author{The Stacks Project Authors}
\date{stacks-project:more-algebra.tex:local-snapshot-20261012; GNU FDL 1.2 or later}
\begin{document}
\maketitle
\setcounter{subsection}{1}
\noindent\textit{Editorial note.} All delivered mathematical text is the author's own, copied from the pinned Stacks Project chapter file more-algebra.tex; no statement or proof step was rewritten, reordered or removed.\par\medskip
\noindent\textit{Editorial note.} The chapter file pinned-original-source.tex carries only \textbackslash{}input\{preamble\} on its first line, so no preamble text lives inside the chapter. The retained file sources/stacks-more-algebra/retained-preamble.tex is a mechanical, line-for-line copy of the project's single shared preamble as staged in the snapshot; nothing was added, removed or redefined and no package was added.\par\medskip
\noindent\textit{Editorial note.} The source URL is the source-level Stacks Project URL, as required when more than one card is produced from one source; the per-claim tag is carried by claim\_key.\par\medskip
\noindent\textit{Editorial note.} The author writes 'by adjoint an nth root' for 'by adjoining', and writes 'b a unit of A'' where the intended reading is that u is the unit; both are reproduced unchanged. The two references inside the counted proof span are recorded as non-load-bearing navigation pointers.\par\medskip
\noindent\textit{Allowed standard prerequisites (editorial).}\begin{itemize}
\item Discrete valuation rings, uniformizers, the residue field and the ramification index. the setting at statement lines 33955-33956 and used throughout the counted proof span. Fixes the objects and the numerical invariants being computed.
\item Normality of A' = A[x]/(x\textasciicircum{}n - pi) and its identification with the integral closure. used at proof lines 33978-33979. Is what upgrades the computation in A' into the structural statement about A\_1.
\item Kummer theory: an extension generated by an n-th root is generated by roots of a divisor when the n-th roots of unity involved lie in the base. invoked at proof lines 33985-33987. Reduces the last statement about subextensions to the unit-root claim that the congruence argument settles.
\item The direct sum decomposition A' = A oplus ... oplus A(pi')\textasciicircum{}\{n-1\} and the shape of the image of multiplication by a power of pi'. used at proof lines 33997-34002. Is the structure that makes the congruence impossible, closing the argument.
\end{itemize}\medskip
\ifdefined\hypertarget\hypertarget{sourcepointer-1}{}\fi
\noindent\textit{Source reference.} Original definition 33041: sources/stacks-more-algebra/pinned-original-source.tex, source lines 33041--33041. Same-chapter navigation pointer; the statement it names is standard for the section and the argument that consumes it stays inside the counted proof span.\par\medskip
\ifdefined\hypertarget\hypertarget{sourcepointer-2}{}\fi
\noindent\textit{Source reference.} Original lemma 0EXP: focused-seeds/AGP/sources/stacks-snapshot-20261011/tags.txt, source lines 18664--18664. Cross-chapter navigation pointer into the Stacks Project; the standard fact it supplies is applied inside the counted proof span and the argument resting on it stays there.\par\medskip
\begin{lemma}
\label{lemma-pull-root-uniformizer}
Let $A$ be a discrete valuation ring with uniformizer $\pi$. Let $n \geq 2$.
Let $K_1 = K[\pi^{1/n}]$. Then
\begin{enumerate}
\item $[K_1 : K] = n$,
\item the integral closure $A_1$ of $A$ in $K_1$ is the ring $A[\pi^{1/n}]$,
\item $A_1$ is a discrete valuation ring,
\item the ramification index of $A_1$ over $A$ is $n$,
\item $K_1$ is totally ramified with respect to $A$, and
\item if $n$ is prime to the residue characteristic of $A$, then
$K_1/K$ is tamely ramified and any subextension of $K_1/K$ is
generated by $\pi^{1/d}$ for some divisor $d$ of $n$.
\end{enumerate}
\end{lemma}

\begin{proof}
Consider the ring $A' = A[x]/(x^n - \pi)$ and denote $\pi'$ the image
of $x$ in $A_1$. As $A'$ is finite free of rank $n$ as an $A$-module,
the element $\pi$ is a nonzerodivisor in $A'$, and hence so is $\pi'$.
Moreover, $A'/\pi'A'$ is isomorphic to $A/\pi A$ which is a field.
Hence $A'$ is a discrete valuation ring with uniformizer $\pi'$.
It follows that the fraction field $K'$ of $A'$ is an extension of degree $n$.
Clearly, $K'$ is obtained from $K$ by adjoint an $n$th root
$\pi' = \pi^{1/n}$ of $\pi$, i.e., $K' \cong K_1$.
Since $A'$ is normal, we see that $A'$ is the
integral closure of $A$ in $K'$, i.e., $A' = A_1$. This proves (1) -- (5).

\medskip\noindent
The assertion that $K_1/K$ is tamely ramified if $n$ is invertible in
the residue field is immediate from
Definition \ref{definition-types-of-extensions} and (3) and (4).
To prove the very last statement, it suffices to show that
every $n$th root of unity $\zeta \in K'$ is contained in $K$, see
Fields, Lemma \ref{fields-lemma-subfields-kummer}.
It is clear that $\zeta \in (A')^*$. We can write $\zeta = a + b$
with $a \in A$ and $b \in A \pi' \oplus \ldots \oplus A (\pi')^{n - 1}$.
If $b$ is nonzero, then we can write $b = u(\pi')^t$ where
$t \geq 1$ and $b$ a unit of $A'$. Then
$$
1 = \zeta^n \equiv a^n + n a^{n - 1} u (\pi')^t \bmod (\pi')^{t + 1}
$$
Since $n$ and $a$ are units of $A$, this would imply that
$b = u(\pi')^t$ is congruent to $(1 - a^n)/(n a^{n - 1}) \in A$
modulo $(\pi')^{t + 1}$. The direct sum decomposition
$A' = A \oplus \ldots \oplus A(\pi')^{n - 1}$ shows that
this is impossible because the image of multiplication by
$(\pi')^{t + 1}$ is a corresponding direct sum whose intersection
with $A(\pi')^t$ is $(\pi A)(\pi')^t$. Thus $b = 0$ and $\zeta \in A$
as desired.
\end{proof}

\end{document}
